\renewcommand{\thetable}{3.A.8}
\begin{tablalafrox}{Restricciones del caso. Fuente: elaboración propia a partir de las restricciones operacionales de Distribuidora Puelche y del Anexo 2.2 del Subdocumento 2 (LafroX, 2026a).}%
  {tab:3-A-8}%
  {F{0.060} Y F{0.124}}%
  {\cab{ID} & \cab{Restricción} & \cab{Origen}}
  R-01 & El despacho de la mañana no se detiene. Entre las 05:30 y las 07:00 salen 96 camiones. No existe forma de ejecutar esa ventana a mano. Un sistema indisponible en ese tramo equivale a un día sin operación. &  \\
  R-02 & El centro de distribución debe poder recibir, preparar y despachar durante un corte del enlace. La preparación nocturna y la carga no pueden quedar detenidas por conectividad. &  \\
  R-03 & Las aplicaciones de preventa y de reparto deben funcionar un turno completo sin señal. Hay rutas donde el dispositivo no recupera cobertura hasta el regreso. &  \\
  R-04 & El sistema de gestión no se reemplaza ni se modifica. Es el sistema de registro contable y tributario. La solución le entrega datos; el sistema de gestión emite los documentos tributarios. No habrá dos verdades ni un segundo emisor. &  \\
  R-05 & No se puede exigir al cliente del canal tradicional que tenga internet, dispositivo propio ni medio de pago electrónico. La solución debe funcionar con el cliente tal como es. &  \\
  R-06 & La solución debe operar con conductores que no son trabajadores de la compañía: 54 camiones de diez empresas transportistas, con rotación sin aviso. &  \\
  R-07 & Los dispositivos que se usen en la cámara de congelado deben operar a -22\,°C y sin cobertura de señal en su interior. &  \\
  R-08 & Congelamiento operacional: prohibido intervenir sistemas entre el 1 y el 25 de septiembre y durante todo el mes de diciembre. Prohibido intervenir durante los tres primeros días hábiles del mes, por el cierre comercial y contable. &  \\
  R-09 & El área de tecnologías de información de la compañía son cuatro personas. Toda función que requiera un especialista dedicado que la compañía no tiene debe ofrecerse como servicio y estar costeada. &  \\
  R-10 & El sindicato de conductores objetó formalmente, el año pasado, la instalación de cámaras en cabina y el control de jornada por posicionamiento satelital. Cualquier propuesta que involucre esas tecnologías debe hacerse cargo del asunto. &  \\
  R-11 & La rotación anual del personal de preparación de pedidos es del 38 \% y el turno crítico es nocturno. Toda solución que exija entrenamiento prolongado para operar en bodega será rechazada. &  \\
  R-12 & La emisión de documentos tributarios electrónicos debe cumplir la normativa vigente. La guía de despacho electrónica y su acuse de recibo tienen efectos legales que la solución no puede comprometer. &  \\
\end{tablalafrox}
